\section{Evaluation}
In this section, we will discuss the basic recipe for "evaluating" in different periods of software development,
starting from the early days of software development, to the current era of Software 2.0 and 3.0, focusing on the differences 
in the evaluation process and the challenges that arise in each era.

\subsection{Software 1.0 - Evaluation in Traditional Software Development}
In traditional software development, the evaluation process is relatively straightforward.
The expected output of the software is well-defined, and we can easily compare the actual output with the expected output to determine if the software is working correctly.
This process is often referred to as "testing", and it can be done through:
\begin{itemize}
    \item \textbf{Unit Testing:} testing individual components or functions of the software to ensure they work as expected.
    For example, if we have a function that adds two numbers, we can test it by providing different pairs of numbers and checking if the output is correct;
    \item \textbf{Integration Testing:} testing the interaction between different components of the software to ensure they work together correctly;
    \item \textbf{Regression Testing:} testing the software to ensure new changes do not break existing functionality.
\end{itemize}
The process is \textbf{deterministic} and \textbf{binary}, meaning that the software either works or it doesn't, and the output is always the same for the same input.
The only difficulty in this process comes from the coverage of the tests, not from the evaluation itself.

\subsection{Software 2.0 - Evaluation in Machine Learning}
In the era of Software 2.0, where machine learning models are used to perform tasks, the evaluation process becomes more complex.
This is because the output of the model is \textbf{not deterministic} anymore, but rather \textbf{probabilistic}, 
The evaluation process, also called \textbf{experimentation}, is no longer binary, but rather a matter of degree. It shifts from "is the output correct?" to "how good is the model across a test set?".

For our analysis, we will focus on the evaluation of a binary classification model, 
where the goal is to classify inputs into two categories (for example, "spam" vs "not spam", or "positive" vs "negative").
Usually the metrics used is the proportion of misclassified inputs (also known as \textbf{error}),
meanwhile the proportion of correctly classified inputs is called \textbf{accuracy}.

\subsubsection{Comparing Models}
Before putting a model into production, we want to compare it with 
other models to understand which one performs better on a given task.
We could also compare different versions of the same model, to understand if the changes we made have improved the performance or not.

In simple cases, such as in 2D dataset, we can plot the predicted output and visually inspect the decision boundary to answer these questions.
However, in more complex cases, such as in high-dimensional datasets, we need to rely on quantitative metrics to compare the models.

The most straightforward approach is to train both models on the same training set, and pick the one with fewer mistakes.
This generates a few questions:
\begin{itemize}
    \item On which data we compute the error? 
    \item How we eliminate random effects? 
    \item Are the metrics we are using good?
\end{itemize}

\subsubsection{Data Splitting}
As you have surely heard before, you should \textbf{never judge your performance on the training set}, 
since the model has already seen them and can easily memorize them, leading to an overestimation of the model's performance on unseen data.

\subsubsection*{Test Set}
The first thing that should be always done in machine learning is to split the original dataset 
into a \textbf{training set} and a \textbf{test set}.
The training set is used to train the model, while the test set is used to evaluate the model's performance on unseen data,
which is more representative of the real-world performance of the model.

\paragraph{How should we split the dataset?}
The only factor to consider when splitting the dataset is the \textbf{size} of the \textbf{test set}. The bigger
it is, the more it is representative of the real-world data distribution, and the more reliable the evaluation results will be.

Ideally, we want 10\,000 examples in the test set, but this is not always possible, especially if the original dataset is small.

\subsubsection*{Overfitting the test set}
However, this approach can lead to overfitting on some cases.
Suppose we want to choose the best hyperparameters for our model, and we are provided with a test set.
We try different hyperparameter settings, and we pick the one that performs best on the test set, supposing that it will also generalize well to unseen data.

At a certain point, we are provided with a fresh test set, and we find out that the performance is much worse than expected with the chosen hyperparameters.
This is because we have \textbf{overfitted the test set}, meaning that we have chosen hyperparameters that perform well on the test set, but not on unseen data.
Our method for choosing the best hyperparameters are another learning algorithm, and by testing different hyperparameters on a small test set,
we are essentially training this learning algorithm on the test set, leading to overfitting.

We can see this as a form of \texttt{multiple testing} problem in statistics. 
We are testing multiple hypotheses (different hyperparameter settings) on the same test set, 
and we are picking the one that performs best, which is likely to be a false positive.

The simple answer to this problem is to \textbf{never use the same test set more than once}, and to always use a fresh test set for the final evaluation of the model's performance.

\subsubsection*{Validation Set}
Reusing the test set not only makes you choose the wrong model, but also \textbf{inflates the performance measures}, 
making them look better than they actually are.

To avoid this, we need to choose the best best model and/or hyperparameters on a separate set of data, 
extracted from the training set and called the \textbf{validation set}.
Ideally, the validation set is the same size of your test set, but you can make it a little smaller to get more data for training.

Note that this approach alone does not solve the problem of multiple testing, but it just provides 
a final fail-safe to detect it,
You need always to be careful and whenever it's needed to rerun your experiments on training and validation sets, 
as you can easily end up overfitting the validation set as well.

However, if your dataset is large, and you haven't done anything strange in tuning phase, 
the validation set results should be a good estimate of the test set results, and you can be reasonably confident that the model will perform well on unseen data.

\subsubsection*{Cross Validation}
In some cases, especially when the dataset is small, we can't afford to set aside a large portion of the data for testing and validation, 
since we need as much data as possible for training.

A technique that can help in this case is called \textbf{Cross validation}.
It allows us to use the entire dataset for both training and testing, by splitting it into multiple folds.
For example, in a 5-fold cross validation, we split the dataset into 5 equal parts, and we train the model on 4 parts and test it on the remaining part (validation).
We repeat this process 5 times, each time using a different part as the test set,and we average the results to get a more reliable estimate of the model's performance.

This can be time-consuming, but ensures that every instance in the dataset is used for training.
After chosen your model and hyperparameters through the cross-validation process, you can test it once on the test set to get an estimate of its performance on unseen data.

\subsubsection*{A special case: Temporal Data}
If your data has a temporal ordering, you need to take it into account when splitting the dataset.
In such cases, you can't sample your test set randomly, since training on future data and testing on past data is not a realistic scenario.
The ordering in the dataset needs to be maintained when splitting the dataset, 
to ensure that the model is trained on past data and tested on future data, which is more representative of the real-world scenario.

However, this is not always the case. If your data has a timestamp, but there's no meaning
associated with the ordering of the data (such as in email classification), you can still split the dataset randomly, since the ordering does not provide any useful information for the task at hand.

If you want to apply cross validation to temporal data, that takes a small portion of initial data as the training set and then iteratively adds more data to the training set, while validating on the next portion of data, until you reach the end of the dataset.

\paragraph{Final remark:}
Philosophically talking, when in doubt make sure that the evaluation setting accurately simulates production, 
and that the validation setting accurately simulates the evaluation setting,

\subsubsection{Experiment Design}
Now that we now how to experiment, what experiments should we run? Which hyperparameters should we test? 
Basically as long as you are not looking at the test set, you can do whatever to like.

\subsubsection*{Hyperparameter Tuning}
When tuning hyperparameters, we have the following options:
\begin{itemize}
    \item \textbf{Trial and Error}: This is the most basic approach, where you simply try different combinations of hyperparameters and see which one performs best on the validation set.
    This is driven by intuition and experience, and it can be time-consuming, especially if you have a large number of hyperparameters to tune.
    \item \textbf{Grid Search}: This is an automated approach, where you define a grid of relevant hyperparameter values, and the algorithm trains and evaluates a model for each combination of hyperparameters in the grid.
    This can be computationally expensive, especially if the grid is large, but it ensures that you are exploring a wide range of hyperparameter combinations.
    \item \textbf{Random Search}: This is another automated approach, where you randomly sample combinations of hyperparameters from a defined distribution, and train and evaluate a model for each combination.
    This can be more efficient than grid search, especially if you have a large number of hyperparameters, since it allows you to explore a wider range of combinations.
\end{itemize}

\subsection{Software 2.0 - Statistical Testing}
Machine Learning and statistics are closely related, but significance testing is rarely used in ML experiments 
due to large sample sizes and emphasis on replication over statistical significance.
Replication is key, but often challenging in practice, so we can still use statistical testing as a tool to assess the reliability of our results.

\subsubsection{True metrics vs Estimated metrics}
When we evaluate a model, we usually imagine data being sampled from a probability distribution $p(x)$.
What we really want is to develop a classifier that performs well on any data point sampled from the original distribution.

Imagine for example we sample one data point from the distribution, and we classify it with some classifier $c$.
The probability of $c$ making a correct prediction on this data point is the \textbf{true accuracy}.
However, we can't compute the true accuracy, since we don't have access to the entire distribution, but only to a finite sample of data points.

Instead, what we usually do is to compute the so-called \textbf{sample accuracy}. 
We take a large enough sample of data points from the distribution, and we compute the accuracy of the classifier on this sample, which is an estimate of the true accuracy.

Now, given the estimated sample accuracy, the question is: 
how sure can we be that the true accuracy is close to the sample accuracy?

\subsubsection*{Confidence Intervals}
One way to quantify the uncertainty around our estimate is to compute confidence intervals, which are ranges of values that are likely 
to contain the true metric with a certain level of confidence.
This helps us understand the uncertainty around our estimate and make more informed decisions about model performance.

The confidence interval is computed based on the variability of the sample metric across different samples, and it gives us a range of values within which we can be reasonably confident that the true metric lies.

The size of the confidence interval depends on two factors:
\begin{itemize}
    \item \textbf{True metric}: If the true metric is close to 0 or 1, the confidence interval will be narrower, since there is less variability in the sample metric.
    \item \textbf{Sample size}: The larger the sample size, the narrower the confidence interval, since there is less variability in the sample metric across different samples.
\end{itemize}
If you have not the luxury of a large sample size, you can do some statistical testing to understand how much the sample metric is close to the true metric.

\paragraph{Remark:} Since we are working with estimated values, the confidence interval itself is an estimate. 
It is not guaranteed to contain the true metric with the specified confidence level for any single realization. 
The interval is relative to the specific test sample collected; therefore, it is not a fixed range but one that varies across different samples.

\subsubsection*{Standard Deviation and Standard Error}
The \textbf{standard deviation} ($\sigma$) measures the amount of variation or dispersion 
within a set of data values. 
In the context of a sampling distribution, the region spanning approximately two standard errors to the left 
and two to the right of the sample mean (a total range of four standard errors) 
contains roughly 95\% of the sample metrics across different samples.
Consequently, this interval $[\bar{x} - 2SE, \bar{x} + 2SE]$ represents a 95\% \textbf{confidence interval} for the true population metric.

The \textbf{standard error} (SE) is the standard deviation of the sampling distribution of a statistic, 
most commonly the sample mean. It quantifies the uncertainty of the sample estimate. It is computed as:
\[
    SEM = \frac{\sigma}{\sqrt{n}}
\]
Where $\sigma$ is the standard deviation of the individual observations and $n$ is the sample size.

While the standard deviation is a measure of the \textbf{spread} of the underlying data, 
the standard error and the confidence interval are measures of \textbf{uncertainty} 
regarding the true population parameter. 
Increasing the sample size $n$ will typically decrease the standard error, thereby narrowing the confidence interval, even if the underlying standard deviation of the population remains constant.


\subsection{Randomness}
\subsection{Randomness}
In machine learning, the results of an experiment can vary significantly upon rerunning due to \textbf{random initialization} 
of model weights and stochastic optimization processes. 
Additionally, \textbf{data sampling} introduces randomness because the available dataset is merely a single realization of a larger, underlying population.

To evaluate robustness to random initialization, executing the experiment multiple times with different seeds is often sufficient. 
To test robustness to data sampling, we typically employ the approaches described in the following sections.

\begin{comment}
    #TODO: go on from this
\end{comment}
\subsubsection{Resampling Methods}
The basic approach is to analyze the \textbf{standard deviation} of the sample metric across different samples, and to compute confidence intervals for the true metric.
Through cross-validation, we can compute the sample metric on different subsets of the data, and analyze how much it varies across these subsets.
This allows us to assess the robustness of our results to data sampling, and to understand how much the sample metric is likely to vary if we were to collect a different sample of data.
\\\\
However, as we increase the number of repeats, we are incrementing the overlap between the samples, especially if we are working with a small dataset.

\subsubsection{Bootstraping Methods}
An alternative approach is to use bootstraping methods, in which we \textbf{resample with replacement} from the original dataset to create multiple bootstrap samples.
This allows us to sample a dataset that is exactly the same size as the original dataset, permitting us to resample as many times as we want without worrying about overlap between samples.
\\
This method approximates the data distribution in a very precisely way.

\subsection{Evaluating Regression}
When computing a continuous metric, such as the mean squared error $m$, we assume individual observations are distributed around a true metric.
The sample mean (computed from the results of the testing) serves as a point estimate of this expectation, but it is itself a random variable.
Across different samples, the sample mean follows a Student's t-distribution which converges to a normal distribution for sample sizes $n > 30$, centered at the true mean.
This allows us to compute confidence intervals for the true metric based on the sample mean and its variability.

\subsubsection{Metrics}
The most common metrics for regression is the mean squared error (MSE), which measures the average squared difference between the predicted values and the true values.
However, also RMSE (root mean squared error) is used, which is the square root of MSE, and it is more interpretable since it is in the same units as the original data, instead of being in squared units.

\subsubsection{Bias and Variance}
Bias and variance are two sources of error in regression models.
Those definitions are strictly related to the MSE, but they can be extended to other metrics as well.
\begin{itemize}
    \item \textbf{Bias} It's the difference between your optimal MSE and the true MSE, and it is a measure of how much your model is underfitting the data.
    This part remains the same also if you resample the data, since it is a property of the model itself, and not of the data.
    \item \textbf{Variance} It's the difference between the true MSE and the measure (across samples), and it is a measure of how much your model is overfitting the data.
    This part can vary across different samples, since it is a property of the data, and not of the model itself.
\end{itemize}
Often the bias and variance are in a trade-off, since increasing the complexity of the model can reduce the bias but increase the variance, and vice versa.

\subsection{Evaluating Classification}
When evaluating classification models, we often use metrics such as accuracy, precision, recall, and F1-score.
These metrics are based on the confusion matrix, which counts the number of true positives, true negatives, false positives, and false negatives.

\subsubsection{Class Imbalance}
It can happen that the classes in the dataset are imbalanced, meaning that one class is much more frequent than the other.
In this case, getting a accuracy of 90\% can be \textbf{misleading}, since the model can simply predict the majority class and still achieve a high accuracy.
\\
In production, this can lead to a model that performs poorly on the minority class, which can be critical in certain applications, such as fraud detection or medical diagnosis.
\\\\
Another problematic case, could be the one where all the minority class is contained in the test set, and the model is not able to learn anything about it, since it is not present in the training set.
This can lead to a model that performs very poorly on the test set, and it is not able to generalize to new data.

\subsubsection{Cost Imbalance}
In some cases, the cost of misclassifying one class can be much higher than the cost of misclassifying the other class.
For example, in medical diagnosis, misclassifying a sick patient as healthy can have much more severe consequences than misclassifying a healthy patient as sick.
\\
In these cases, it's important to use metrics that take into account the class imbalance and cost imbalance.
We can tune our error by multiplying the error of each class by a cost factor, which can be determined based on the specific application and the relative importance of the classes.
\\\\
It's important to note that, in both cases, the cost is never zero, even in the example of medical diagnosis, since misclassifying a healthy patient as sick can lead to unnecessary treatments and anxiety for the patient.
Especially in social applications, where could bring to a lot of ethical issues.

\subsubsection{Confusion Matrix}
The confusion matrix is a table that summarizes the performance of a classification model by counting the number of true positives, true negatives, false positives, and false negatives.
This allows us to compute various metrics.
\\\\
However it's difficult to compare two models based on the confusion matrix, but it can give good insights on the single model works.
The problems previously described affects the diagonal of the confusion matrix (false positives and false negatives - misclassifications) %He wrotes this, but doesn't make sense
\\\\
You can plot the confusion matrix both for the training set, validation set, and test set, giving different insights on the three sets.

\subsubsection{Precision and Recall}
\textbf{Precision} is the ratio of true positives to the total number of predicted positives, while \textbf{recall} is the ratio of true positives to the total number of actual positives.
In order to handle class imbalance, we need to trade-off between precision and recall, since increasing precision can lead to a decrease in recall, and vice versa.
\\
To help us in this trade-off, we can plot the precision and recall for different classifier.

\subsubsection{TPR and FPR - ROC space}
\textbf{TPR} (True Positive Rate) is the ratio of true positives to the total number of actual positives, while \textbf{FPR} (False Positive Rate) is the ratio of false positives to the total number of actual negatives.
We want to maximize TPR and minimize FPR, but as before there is a trade-off between the two.
The optimal space is TPR/FPR = 1, which means that we are maximizing TPR and minimizing FPR at the same time, this is called the \textbf{ROC space} (Receiver Operating Characteristic).

\subsubsection{Plotting thesholds}
Until now we have been talking about the comparison of different models, but we can also compare the same model with different thresholds (eager to predict positive).
If we have a shy model (high threshold), we would not predict anything, however if we have a bold model (low threshold), we would predict everything as positive even risking misclassifications.
The curve that we have, it can be a good way to compare the performance of the model across different thresholds, and to choose the optimal threshold based on the specific application and the trade-off between precision and recall.
\\\\
However, this is possible only if the model outputs is \textbf{ranking classifier} (i.e., it provides a score for each instance), which is not always the case, since some models output a binary prediction without a score.
In this case, the \textbf{ranking error} is the error measured in confidence, if we predict a positive class with a high confidence, but it is actually a negative class, this would be an high ranking error.

\subsubsection{Coverage Matrix}
Given a model that outputs a ranking/score, we can plot the coverage matrix, with negative class on the x-axis and positive class on the y-axis.
Basically the more confident the model is predicting a positive class, the closer the point is to the top right corner, while the more confident the model is predicting a negative class, the closer the point is to the bottom left corner.
Then coloring binary the points based on the predicted class (after thesholding), we can have a good visual representation of the performance of the model.

\subsubsection{AUC}
The AUC (Area Under the Curve) is a metric that summarizes the performance of a classification model across all possible thresholds.
It is computed as the area under the ROC curve, and can be used to compare different models.
If we normalize the coverage matrix, we get the ROC curve, and the AUC is the area under this curve.
The AUC (in ROC space) is an estimate of the probability that a ranking classifier puts a randomly drawn pair of positive and negative examples in the correct order.
\\\\
Lets now look how to works for another type of classifier.
The decision tree is an example of a \textbf{partitioning classifier}, and how we get a ranking from a classifier that depends entirely on the model class, in other words, it depends on the way the model is built.
\\
If we choose to use the proportion of positive point for each segment (leaf) as the score, and we then put all the points in one region on the same level in the ranking.
This means that some pairs of points will have the same score, and this can lead to ties in the ranking.
Then, for large datasets, these regions will not contribute much to the total curve.

\subsubsection{Precision-Recall Curve}
The Precision-Recall curve is another way to evaluate the performance of a classification model, similar to the ROC curve, but in the axis we have precision and recall instead of TPR and FPR.
The area under the Precision-Recall curve (AUC-PR) can be used as a metric to compare different models, especially in cases of class imbalance, where the ROC curve can be misleading.
\\
However in some cases, the Precision-Recall curve can be more difficult to interpret than the ROC curve, since it is not normalized and it can be affected by the class imbalance in a different way than the ROC curve.
Both curves can be useful in different scenarios, it's suggested to plot both.

\subsubsection{Use of Metrics}
There are many others metrics that can be used to evaluate classification models, and they can be applied to all the sets (training, validation, and test set).
For instance, if we compute the accuracy on the test set, we are talking about the \textbf{test accuracy}, used to show if our conclusions are valid.
Meanwhile, the \textbf{validation accuracy} is used to tune the hyperparameters of the model.
And the \textbf{training accuracy}, which is used in comparison with the validation accuracy (difference \textbf{generalization gap}) to understand if the model is overfitting or underfitting the data.
It's important to note that it's not used for evaluating on the training set.

\subsection{No Free Lunch Theorem}
A question that often aries is: which classifier, model, search methods is the best?
The aim is to make a guess before seeing data, but the No Free Lunch Theorem states that is not possible.
\\\\
If you average over every possible problem, no optimization or learning algorithm is better than any other.
An algorithm only looks good when the problem happens to match its assumptions.
This means that there is no universally best model or algorithm. (\textbf{problem of induction})
\\\\
The \textbf{inductive bias} of a learning algorithm or model are those assumptions about the domain that are explicitly or implicitly hardcoded into the model.
In other words, the inductive bias is the set of assumptions that a model makes about the data in order to generalize from the training data to unseen data.
As a computer scientist, we need to figure out which of the many possible inductive biases is the best for our problem.
For instance a linear model has a strong inductive bias towards linear relationships between features and the target variable.
\\\\
Finally, when two models explain the data equally well, we usually prefer the simpler one.
The intuition is that simple patterns are more likely than very complex ones, so learning methods often include a built-in \textbf{bias toward simplicity}.